% Imported from LaTeX_and_pdfs/DGDT_dump.tex; see provenance.json.
\part{Analytical Geometry}


\section{Applications of Vector Algebra to Analytic Geometry}

cf. Ch. 13 "Applications of Vector Algebra to Analytic Geometry", Apostol (1991) \cite{dgdt:Apos1991}

\subsection{Lines in $n$-space}

cf. Sec. 13.2 "Lines in $n$-space", Apostol (1991) \cite{dgdt:Apos1991}.

\begin{definition}[Line]
	Let given pt. $P$, given $A$ nonzero vector. 
	
	Define line $L(P;A)$
	\begin{equation}
		L(P;A) := \lbrace P + tA | t\in \mathbb{R} \rbrace \equiv L(P;A) = \lbrace P + tA \rbrace
	\end{equation}
\end{definition}

\subsubsection{Some simple properties of a straight line}

cf. Sec. 13.3 "Some simple properties of straight line", Apostol (1991) \cite{dgdt:Apos1991}.

\begin{theorem}[Thm. 13.4 of Apostol (1991) \cite{dgdt:Apos1991}]
	2 distinct pts. determine a line, i.e. if $P\neq Q$, $\exists !$ line containing both $P$ and $Q$, namely
	\begin{equation}
		\lbrace P + t(P- Q) \rbrace
	\end{equation}
\end{theorem}

\begin{proof}
Let $L$ be line parallel through $P$ to $Q-P$, i.e. $L= L(P;Q-P)= \lbrace P + t(Q-P) \rbrace$.

$P,Q \in L$ for $t=0,1$, respectively.

Suppose $L'$ is another line containing $P$ and $Q$. Since $P \in L'$, $L' = L'(P,A)$, $Q \in L'$ so $Q = P + At$. \\
$Q - P = tA$. $t\neq 0$ since $Q \neq P$. So $A = \frac{Q- P}{t}$, i.e. $A$ and $Q-P$ are parallel. \\

By Thm. 13.1, since $L(P;Q-P), L'(P,A)$ through same pt. $P$ and direction vectors $Q-P,A$ parallel, then $L=L'$.

\end{proof}


\subsection{Planes in Euclidean $n$-plane}

cf. Sec. 13.6 "Planes in Euclidean $n$-space", Apostol (1991) \cite{dgdt:Apos1991}.

\begin{definition}
	Set $M$ of pts. in $V_n$ is called a plane if $\exists$ pt. $P$ and 2 linearly independent vectors $A$ and $B$ s.t.
	\begin{equation}
		M = \lbrace P + sA + tB | s, t \in \mathbb{R} \rbrace
	\end{equation}
\end{definition}

EY 20211207: Consider the intersection of 2 planes: $M \cap N$:
\[
\begin{aligned}
	& M = \lbrace P + sA + tB | s, t \in \mathbb{R} \rbrace \\
	& N = \lbrace Q + uC + vD \rbrace
\end{aligned}
\]

Let $X\in M$ and $X\in N$ so $X = P+sA + tB = Q+ uC + vD$. 

Let $M\cap N \neq \emptyset$. Let $X_0 \in M \cap N$, $X_0 = P + s_0 A + t_0 B = Q + u_0 C + v_0 D$.


